\chapter{Cross-Currency Basis and FX Carry}\label{ch:s2:cross-currency-basis-and-fx-carry}

A basket long the three G10 currencies with the highest interest rates against the dollar and short the three lowest earned 2.7\% a year from May
2002 to December 2025, built here from public exchange rates and interbank rates. Its Sharpe ratio was 0.37, and it lost 19.4\% from September to
November 2008. Jurek hedged such baskets with out-of-the-money options and found that crash risk premia account for at most a third of the carry trade's
return. The other currency trade of this chapter is the cross-currency basis. Du, Tepper and Verdelhan found that borrowing dollars through FX swaps
has cost more than covered parity says, persistently, and most for contracts that sit on banks' balance sheets at quarter-ends. On this chapter's
synthetic market, a put-hedged carry basket gives up 40\% of its return to lose its negative skew, and lending dollars through the swap earns 11.2
basis points a year after a \omterm{def:s2:swap-spread-and-asset-swap-trades:bscost}{balance-sheet cost}. The build is \texttt{firm.fxcarry}.

\section{FX carry: construction and crash risk}

\begin{definition}[FX carry basket]\label{def:s2:cross-currency-basis-and-fx-carry:basket}
An \emph{FX carry basket}\index{FX carry basket} holds long the currencies with the highest short-term interest rates and short those with the lowest,
usually in equal weights and against a common base currency, rebalanced as rates change; its return is the rate differential earned plus the change
in the exchange rates.
\end{definition}

Book~8 (chapter~20) built the carry strategy across asset classes and its crashes. Here the currency version is built from public data. Each month,
on the previous month's three-month interbank rates, the basket ranks nine currencies (euro, yen, sterling, Swiss franc, Australian and New Zealand
dollars, Canadian dollar, Swedish and Norwegian kronor) by their rate minus the US rate. It goes long the top three and short the bottom three against
the dollar. From May 2002 to December 2025 (284 months) it earned 2.7\% a year with a volatility of 7.2\%, a Sharpe ratio of 0.37 and a skewness of
$-0.60$ (\cref{fig:s2:cross-currency-basis-and-fx-carry:g10}). The rate differential earned 3.1\% a year and the exchange rates took back 0.4\%. Its
worst month was October 2008, $-10.5\%$. The New Zealand dollar was long in 271 of the 284 months, the Australian dollar in 218, the krone in 188. The
Swiss franc was short in all 284, the yen in 177, the krona in 166, the euro in 159.

\begin{omfigure}
\begin{tikzpicture}
\begin{axis}[width=11cm,height=5.4cm,xlabel={\small year},ylabel={\small cumulative log return (\%)},tick label style={font=\footnotesize},
  xmin=2002,xmax=2026.5,ymin=0,ymax=65,grid=major,grid style={gray!25},table/col sep=comma,/pgf/number format/1000 sep={}]
\addplot[omThm,thick] table[x=year,y=cum]{figdata/strategies-2/14-cross-currency-basis-and-fx-carry/g10.csv};
\end{axis}
\end{tikzpicture}

\omcaption{A G10 carry basket against the dollar (long the three highest three-month interbank rates, short the three lowest), cumulative log return,
May 2002 to December 2025, from FRED exchange rates and OECD interbank rates. Derived statistics only. Data:
\texttt{s2\_fetch\_g10}.}\label{fig:s2:cross-currency-basis-and-fx-carry:g10}
\end{omfigure}

The public-data basket is approximate. The OECD's rates are monthly and differ by country in how they are measured. A trader would use forwards, whose
points give the rate differential exactly, and would pay their spreads. The shape is the familiar one: steady gains, then a crash when high-yielding
currencies fall together.

\subsection*{Crash hedges}

\begin{definition}[Crash-hedged carry]\label{def:s2:cross-currency-basis-and-fx-carry:hedged}
\emph{Crash-hedged carry}\index{crash-hedged carry} is a carry position combined with out-of-the-money options that pay when the high-yielding
currencies fall sharply against the funding currencies, so that the position's return no longer carries the crash; the options' cost measures the
price of the crash risk the unhedged trade bears.
\end{definition}

Jurek used G10 option prices from 1990 to 2012. Carry trades earned Sharpe ratios as high as equity market factors or higher, and hedging with
out-of-the-money options showed that the returns were not a peso problem. Crash risk premia accounted for at most a third of them.

\texttt{firm.fxcarry} runs the same comparison on the ten synthetic currencies of Book~8's futures universe, which carry a premium and three planted
crashes (\cref{lst:s2:cross-currency-basis-and-fx-carry:hedge}). The basket earned 1.30\% a year with a Sharpe ratio of 0.23, a skewness of $-0.39$ and
a worst month of $-6.6\%$. A monthly put on the basket, struck 1.5 standard deviations below zero and priced at 1.2 times the trailing volatility, cost
1.14\% a year and paid 0.63\%. Hedged, the basket earned 0.78\% a year, a Sharpe ratio of 0.15, with no skew (0.01) and a worst month of $-4.3\%$
(\cref{fig:s2:cross-currency-basis-and-fx-carry:hedge}). The planted crashes lasted a hundred days each and took 14.1\%, 15.5\% and 14.1\% from the
basket. Over the months they spanned, the hedge cut the first from $-12.0\%$ to $-7.9\%$ and the third from $-14.4\%$ to $-8.5\%$. It did almost
nothing for the second ($-14.7\%$ to $-14.3\%$), which fell a little each month, never far enough below the strike. The hedge's net cost was 40\% of
the carry return, more than Jurek's third. That is a planted choice of option prices.

\begin{omfigure}
\begin{tikzpicture}
\begin{axis}[width=11cm,height=5.4cm,xlabel={\small year},ylabel={\small cumulative log return (\%)},tick label style={font=\footnotesize},
  xmin=0,xmax=30,ymin=-15,ymax=50,grid=major,grid style={gray!25},table/col sep=comma,
  legend style={font=\scriptsize,at={(0.03,0.97)},anchor=north west}]
\addplot[omThm,thick] table[x=year,y=raw]{figdata/strategies-2/14-cross-currency-basis-and-fx-carry/hedge.csv};
\addplot[omDef,thick] table[x=year,y=hedged]{figdata/strategies-2/14-cross-currency-basis-and-fx-carry/hedge.csv};
\legend{carry basket,with monthly puts}
\end{axis}
\end{tikzpicture}

\omcaption{The synthetic currency carry basket over thirty years with three planted crashes, unhedged and with a monthly put struck 1.5 standard
deviations down. Data: \texttt{s2\_fxcarry.cumulative}.}\label{fig:s2:cross-currency-basis-and-fx-carry:hedge}
\end{omfigure}

\section{The cross-currency basis}

Covered interest parity (Book~2, chapter~16) says that borrowing dollars directly and borrowing another currency and swapping it into dollars should
cost the same. The cross-currency basis measures by how much they differ; a negative basis means the swap route costs more, so dollars are dear through
FX swaps. Du, Tepper and Verdelhan found the deviations large, persistent, systematic, and not explained by credit risk or transaction costs. They
correlated with other fixed-income spreads.

\begin{definition}[Basis funding trade]\label{def:s2:cross-currency-basis-and-fx-carry:funding}
A \emph{basis funding trade}\index{basis funding trade} lends dollars through FX swaps (receiving the other currency and paying it back forward) when
the cross-currency basis makes that route pay more than lending dollars directly, earning minus the basis less the cost of the balance sheet the
position uses.
\end{definition}

The synthetic basis has a level of $-20$ basis points a year and noise, plus quarter-end windows: in the last ten trading days of each quarter it is
30 basis points more negative, 60 at year-end. It averaged $-26.2$ basis points over ten years, $-49.8$ in the quarter-end windows and $-80.4$ in the
year-end ones. A lender whose balance sheet costs 15 basis points a year earned 11.2 basis
points a year by lending all the time. Lending only in the windows, 40 days a year, earned 6.7, at a rate of 42.5 basis points a year while lent.

\section{Quarter-ends}

The quarter-end pattern is the balance-sheet argument in its plainest form. Nothing about credit or rates changes on 31 March, but banks that report
their balance sheets on that date shrink them over it, and the dollars they would have lent through the swap become scarce. The chapter's windows are
planted. The real ones are only as large as the balance-sheet rules that cause them. A lender that is not itself measured on those dates is paid to
provide the balance sheet others withdraw.

\section{Funding-driven trades}

The two halves of the chapter meet in funding. A carry trade funds itself in the low-rate currency. When dollars are dear through swaps, the forward
points that set the carry include the basis. A carry position hedged in the swap market pays the basis, and one financed in cash does not. Carry books
also crash when funding tightens. The high-yielders fall as their holders sell to repay funding-currency loans. A momentum filter that cuts the
basket when its own trend turns negative is the usual defence. It gives up part of the carry in the calm and misses the first days of any crash.

\section{Strategy files}

\begin{strategyfile}{G10 carry basket}\label{strat:s2:cross-currency-basis-and-fx-carry:basket}
\sfield{Who pays you, and why} Investors who pay to hold low-yielding, safe currencies; the risk of high-yielders crashing in a flight to safety.
\sfield{Instruments and venues} FX forwards; futures.
\sfield{Signal} Forward points (the rate differential).
\sfield{Sizing and execution} Top three long, bottom three short, equal weights; monthly.
\sfield{Costs} Forward spreads; roll.
\sfield{How it dies} Crashes in which high-yielders fall together: $-19.4\%$ in autumn 2008 for the public-data basket.
\sfield{Horizon, capacity, infrastructure} Months to years; large capacity.
\sfield{Backtest honestly} Forward points, not interbank rates; transaction costs; the 2008 crash in the sample.
\sfield{Sources} Jurek (2014); this chapter's FRED/OECD basket: 2.7\% a year, Sharpe ratio 0.37, 2002--2025.
\end{strategyfile}

\begin{strategyfile}{Crash-hedged carry with options}\label{strat:s2:cross-currency-basis-and-fx-carry:hedged}
\sfield{Who pays you, and why} As for carry, less the crash premium paid to option sellers.
\sfield{Instruments and venues} Carry positions and out-of-the-money FX options on each pair or on the basket.
\sfield{Signal} Carry; the hedge is a programme.
\sfield{Sizing and execution} Monthly options struck one to two standard deviations out.
\sfield{Costs} Option premiums and spreads.
\sfield{How it dies} Slow crashes that never reach the strike.
\sfield{Horizon, capacity, infrastructure} Monthly; option access.
\sfield{Backtest honestly} Option prices from the smile at the time, not from a flat vol.
\sfield{Sources} Jurek (2014): crash premia at most a third of carry returns; this chapter: the hedge cost 40\% and removed the skew.
\end{strategyfile}

\begin{strategyfile}{Quarter-end basis trade}\label{strat:s2:cross-currency-basis-and-fx-carry:qe}
\sfield{Who pays you, and why} Banks that shrink their balance sheets over reporting dates.
\sfield{Instruments and venues} FX swaps spanning the quarter-end.
\sfield{Signal} The calendar; the basis priced for the turn.
\sfield{Sizing and execution} Lend dollars across the turn; limited by one's own balance sheet on the date.
\sfield{Costs} One's own \omterm{def:s2:swap-spread-and-asset-swap-trades:bscost}{balance-sheet cost} across the turn.
\sfield{How it dies} Reporting rules that average over the quarter.
\sfield{Horizon, capacity, infrastructure} Days to weeks.
\sfield{Backtest honestly} Swap prices for the exact turn dates.
\sfield{Sources} Du, Tepper and Verdelhan (2018); this chapter: 42.5 basis points a year while lent in the windows.
\end{strategyfile}

\begin{strategyfile}{Basis funding arbitrage}\label{strat:s2:cross-currency-basis-and-fx-carry:funding}
\sfield{Who pays you, and why} Borrowers of dollars through FX swaps (non-US banks and investors hedging dollar assets) who pay above parity.
\sfield{Instruments and venues} FX swaps and cross-currency swaps; dollar funding.
\sfield{Signal} The basis against one's own funding and \omterm{def:s2:swap-spread-and-asset-swap-trades:bscost}{balance-sheet cost}.
\sfield{Sizing and execution} As balance sheet allows; term matched.
\sfield{Costs} Balance sheet; counterparty lines.
\sfield{How it dies} A funding squeeze on the lender itself.
\sfield{Horizon, capacity, infrastructure} Months; treasury-level funding.
\sfield{Backtest honestly} Charge the lender's real \omterm{def:s2:swap-spread-and-asset-swap-trades:bscost}{balance-sheet cost}.
\sfield{Sources} Du, Tepper and Verdelhan (2018); this chapter: 11.2 basis points a year after a 15 basis-point cost.
\end{strategyfile}

\begin{strategyfile}{Carry with momentum filter}\label{strat:s2:cross-currency-basis-and-fx-carry:momentum}
\sfield{Who pays you, and why} As for carry, with some crash avoidance paid for by lost carry.
\sfield{Instruments and venues} As for carry.
\sfield{Signal} Carry, cut when the basket's trailing return is negative.
\sfield{Sizing and execution} Reduce or close when the filter trips.
\sfield{Costs} More turnover.
\sfield{How it dies} Whipsaws; fast crashes that begin before the filter trips.
\sfield{Horizon, capacity, infrastructure} Months.
\sfield{Backtest honestly} The filter's parameters chosen before the test.
\sfield{Sources} No performance figure verified.
\end{strategyfile}

\section{Tutorial: priced at the quarter-end}\label{tut:s2:cross-currency-basis-and-fx-carry:tut}

\begin{tutorial}
\textbf{Goal.} Build carry baskets, hedge their crashes with options, simulate a cross-currency basis with quarter-end windows and lend into it;
build the G10 basket from public data. \textbf{End state:} the numbers in the text and the two figures.
\begin{tutsteps}
\item \textbf{Basket and hedge}.
\omcode{code/firm/fxcarry/firm_fxcarry.py}{27}{64}{Carry baskets and a monthly put on the basket.}\label{lst:s2:cross-currency-basis-and-fx-carry:hedge}
\item \textbf{Carry results}.
\omcode{code/strategies-2/14-cross-currency-basis-and-fx-carry/python/s2_fxcarry.py}{45}{55}{Unhedged and hedged statistics, and the crash months.}\label{lst:s2:cross-currency-basis-and-fx-carry:results}
\item \textbf{Run} \texttt{s2\_fetch\_g10.py} once, then \texttt{carry\_results()}, \texttt{basis\_results()} and \texttt{fig\_fxcarry.py}.
\end{tutsteps}
\textbf{What to change next.} Hedge each pair with its own option instead of the basket; add a momentum filter and compare the crashes; tie the basis
to the carry basket's funding.
\end{tutorial}

\section{Build: FX carry and basis}\label{bld:s2:cross-currency-basis-and-fx-carry:fxcarry}

\begin{build}
\bfield{Purpose} Carry baskets, option-hedged carry, a cross-currency basis with quarter-end windows, and \omterm{def:s2:cross-currency-basis-and-fx-carry:funding}{basis funding trades}.
\bfield{Interface} \texttt{carry\_basket(r, carry, k, every)}, \texttt{crash\_hedge(basket, every, strike\_sd, vol\_mult, lookback)},
\texttt{BasisConfig(\dots)}, \texttt{simulate\_basis(cfg)}, \texttt{basis\_trade(b, cfg, days)}.
\bfield{Rules} Rebalance on yesterday's carry; puts priced at the start of each period on trailing vol; basis trades earn minus the basis less the
\omterm{def:s2:swap-spread-and-asset-swap-trades:bscost}{balance-sheet cost}.
\bfield{Acceptance tests} \texttt{code/firm/fxcarry/tests/}: the basket's positions by hand; the put's price and payoff by hand; the basis windows
and the trade's daily accrual.
\bfield{Stretch} Pair-level options; forward points with the basis; funding crashes.
\end{build}

\begin{omsources}
\item J.~W.~Jurek, ``Crash-neutral currency carry trades'', \emph{Journal of Financial Economics} 113(3), 2014.
\item W.~Du, A.~Tepper and A.~Verdelhan, ``Deviations from covered interest rate parity'', \emph{Journal of Finance} 73(3), 2018.
\item Board of Governors H.10 exchange rates and OECD interbank rates, via FRED.
\end{omsources}

\section{Exercises}

\begin{exercise}[$\star$]\label{exo:s2:cross-currency-basis-and-fx-carry:1}
The basket earned 2.7\% a year with a volatility of 7.2\%. What is its Sharpe ratio, without a risk-free rate?
\end{exercise}

\begin{exercise}[$\star$]\label{exo:s2:cross-currency-basis-and-fx-carry:2}
The basis is $-50$ basis points a year and a lender's balance sheet costs 15. What does lending for ten trading days earn, in basis points of the
amount?
\end{exercise}

\begin{exercise}[$\star$]\label{exo:s2:cross-currency-basis-and-fx-carry:3}
A hedge costs 1.14\% a year and pays 0.63\%. What is its net cost as a share of a 1.30\% carry return?
\end{exercise}

\begin{exercise}[$\star\star$]\label{exo:s2:cross-currency-basis-and-fx-carry:4}
Why was the Swiss franc short in every month of the public-data basket?
\end{exercise}

\begin{exercise}[$\star\star$]\label{exo:s2:cross-currency-basis-and-fx-carry:5}
Why did the monthly put fail to protect against the second synthetic crash?
\end{exercise}

\begin{exercise}[$\star\star$]\label{exo:s2:cross-currency-basis-and-fx-carry:6}
Why does a negative cross-currency basis persist?
\end{exercise}

\begin{exercise}[$\star\star\star$]\label{exo:s2:cross-currency-basis-and-fx-carry:7}
\emph{Coding.} Run \texttt{carry\_results(vol\_mult=1.0)}, pricing the puts at trailing volatility with no premium. What does the hedge now cost, and
what does that say about where the crash premium sits?
\end{exercise}

\begin{exercise}[$\star\star\star$]\label{exo:s2:cross-currency-basis-and-fx-carry:8}
\emph{Find the flaw.} ``Carry earned 2.7\% a year with a Sharpe ratio of 0.37 over 23 years; we lever it to 20\% volatility for 7.5\% a year.''
\end{exercise}

\section{Problem: Priced at the Quarter-End}

\begin{problem}[{Weekend problem --- currencies and balance sheets}]\label{pb:s2:cross-currency-basis-and-fx-carry:1}
The chapter's synthetic markets, the public-data basket and the public record.

\medskip\noindent\textbf{Part I --- Carry.}
\begin{enumerate}
\item Define an \omterm{def:s2:cross-currency-basis-and-fx-carry:basket}{FX carry basket} and build it from public data.
\item Give its return, risk and crash.
\item Which currencies did it hold, and how often?
\item What did Jurek find?
\end{enumerate}

\medskip\noindent\textbf{Part II --- Hedges.}
\begin{enumerate}[resume]
\item Define \omterm{def:s2:cross-currency-basis-and-fx-carry:hedged}{crash-hedged carry}.
\item Give the synthetic basket's statistics with and without the hedge.
\item What did the hedge do in each crash?
\item Why does the hedge's cost measure the crash premium?
\end{enumerate}

\medskip\noindent\textbf{Part III --- The basis.}
\begin{enumerate}[resume]
\item Define covered interest parity and the cross-currency basis.
\item What did Du, Tepper and Verdelhan find?
\item Define the \omterm{def:s2:cross-currency-basis-and-fx-carry:funding}{basis funding trade} and give its synthetic return.
\item What does the quarter-end window add?
\end{enumerate}

\medskip\noindent\textbf{Part IV --- The verdict.}
\begin{enumerate}[resume]
\item State the \emph{named result}: the carry basket's Sharpe ratio with and without crash hedges, and the basis trade's quarter-end return.
\item How do carry and the basis meet in funding?
\item What would a momentum filter change?
\item How would you size a carry book?
\item How would you backtest the basis trade honestly?
\item Which strategy file needs its own balance sheet?
\item How does this chapter relate to chapter~12?
\item In one sentence: what does the cross-currency basis price?
\end{enumerate}
\end{problem}

\section{Interview questions}

\begin{interviewq}[$\star$ \iqroles{trader}]\label{iq:s2:cross-currency-basis-and-fx-carry:1}
What is the currency carry trade, and why does it work?
\end{interviewq}

\begin{interviewq}[$\star\star$ \iqroles{researcher}]\label{iq:s2:cross-currency-basis-and-fx-carry:2}
How would you measure the crash risk premium in currency carry?
\end{interviewq}

\begin{interviewq}[$\star\star$ \iqroles{trader}]\label{iq:s2:cross-currency-basis-and-fx-carry:3}
The three-month cross-currency basis is $-40$ basis points. What does it mean, and who pays it?
\end{interviewq}

\begin{interviewq}[$\star\star$ \iqroles{risk}]\label{iq:s2:cross-currency-basis-and-fx-carry:4}
How would you stress a G10 carry book?
\end{interviewq}

\begin{interviewq}[$\star\star$ \iqroles{developer}]\label{iq:s2:cross-currency-basis-and-fx-carry:5}
Design the data behind a carry book: rates, forwards, spot, holidays, and what can go wrong at month-end.
\end{interviewq}

\begin{interviewq}[$\star\star\star$ \iqroles{researcher}]\label{iq:s2:cross-currency-basis-and-fx-carry:6}
Derive the forward exchange rate from covered interest parity, and show how a basis $b$ enters it.
\end{interviewq}
